\section{Conclusiones y futuro trabajo}

\subsection{Conclusiones}

El desarrollo de este proyecto permitió evaluar la viabilidad de la predicción temprana del shock hemorrágico postoperatorio en pacientes sometidos a cirugía mayor no cardíaca mediante modelos de inteligencia artificial. A partir de cinco experimentos secuenciales (baseline, agregaciones, categorizaciones, pruning e hiperparámetros), se obtuvieron las siguientes conclusiones:

\begin{enumerate}
    \item \textbf{Desempeño global moderado pero clínicamente utilizable:} El mejor desempeño global se alcanzó en el Experimento 5 con \texttt{random\_forest} (AUC-ROC CV = 0.6404). Frente al baseline del Experimento 1 (AUC-ROC = 0.6153), la mejora absoluta fue de 0.0251.

    \item \textbf{Importancia de la ingeniería de variables:} Las agregaciones del Experimento 2 mejoraron el rendimiento respecto al baseline (0.6192 vs 0.6153), mientras que las categorizaciones del Experimento 3 redujeron el AUC (0.6134). Este patrón respalda la hipótesis de pérdida de información al discretizar variables continuas \cite{ma2023}.

    \item \textbf{Pruning como etapa de depuración útil:} En el Experimento 4 se eliminaron siete variables de baja ocurrencia y aporte marginal, alcanzando el mejor rendimiento intermedio (0.6237), lo cual indica que la reducción de ruido favorece la capacidad discriminativa de los modelos.

    \item \textbf{Variables con mayor señal predictiva:} En el modelo campeón (Experimento 5, \texttt{random\_forest}), las variables con mayor importancia fueron \texttt{EDAD} (52.71\%), \texttt{HB\_PREQX} (27.93\%) y \texttt{SANGRADO\_MAYOR} (15.36\%). Esto sugiere que edad, hemoglobina preoperatoria y sangrado intraoperatorio constituyen los pilares de predicción para shock hemorrágico postoperatorio.

    \item \textbf{Comparación de algoritmos y decisión final:} Se evaluaron siete algoritmos (CatBoost, Decision Tree, LightGBM, Regresión Logística Lasso, Naive Bayes, Random Forest y XGBoost). El mejor AUC fue de \texttt{random\_forest} (0.6404), seguido de \texttt{lightgbm} (0.6384) y \texttt{xgboost} (0.6371). El modelo interpretable \texttt{decision\_tree} se posicionó cuarto con AUC=0.6354, consolidando una alternativa de transparencia cercana al mejor desempeño.

    \item \textbf{Trade-off clínico en el punto operativo:} Bajo priorización de sensibilidad, el modelo \texttt{random\_forest} operó con umbral óptimo de 0.38, recall de 0.8196 y especificidad de 0.3567. Esto implica una tasa de falsos positivos de aproximadamente 64.33\%, por lo que el riesgo de fatiga de alertas debe gestionarse en implementación clínica \cite{steyerberg2018}.

    \item \textbf{Interpretabilidad como requisito de adopción:} Aunque \texttt{random\_forest} fue el modelo de mayor discriminación, \texttt{decision\_tree} ofrece mayor transparencia para discusión clínica. El árbol final completo se documenta en anexos como soporte de trazabilidad del modelo. La brecha de AUC entre ambos es de solo 0.0050, lo cual representa un trade-off razonable entre desempeño e interpretabilidad.
\end{enumerate}

\subsection{Limitaciones del estudio}

\begin{itemize}
    \item \textbf{Definición de la variable objetivo:} La variable objetivo (shock hemorrágico) fue definida como variable booleana binaria (shock presente/ausente) basada en el Índice de Shock (IS = FC/PAS), en lugar de utilizar criterios clínicos precisos que distingan entre clases de shock (compensado, descompensado, irreversible). Esta simplificación limita la capacidad del modelo para capturar la severidad y el estadio evolutivo del shock hemorrágico.

    \item \textbf{Tamaño y cobertura del dataset:} El tamaño muestral y la naturaleza de los datos limitan la capacidad de generalización y el techo de desempeño discriminativo. Como señalan Steyerberg et al.\ \cite{steyerberg2018}, los resultados de evaluación interna no garantizan transferencia directa al entorno clínico real.

    \item \textbf{Desbalance de clases y SMOTE:} La proporción 2.1:1 entre casos negativos y positivos requirió técnicas de balanceo (SMOTE, class weights). El sobremuestreo sintético en datasets pequeños puede crear patrones artificiales, lo cual explica parcialmente los altos valores de recall con baja especificidad observados en algunos experimentos.

    \item \textbf{Manejo de variables continuas:} Los experimentos de categorización de EDAD y HB\_PREQX no produjeron mejoras en el rendimiento. Este resultado está en línea con la revisión sistemática de Ma et al.\ \cite{ma2023}, que demuestra que la dicotomización y categorización de predictores continuos conlleva pérdida de información y reduce la precisión de los modelos de predicción clínica. En nuestro caso, el tratamiento de estas variables como continuas resultó más apropiado.

    \item \textbf{Fuentes de datos unimodales:} A diferencia de enfoques como ShockModes \cite{pal2022}, que combina notas clínicas, series temporales de signos vitales y comorbilidades, este proyecto cuenta únicamente con variables tabulares estáticas. La ausencia de variables dinámicas (signos vitales en tiempo real, tendencias de presión arterial) limita la capacidad predictiva del modelo, ya que el shock hemorrágico es un proceso evolutivo que no se captura completamente con datos estáticos.

    \item \textbf{Validación externa:} El modelo no fue validado en poblaciones externas a la Fundación Valle del Lili, lo que limita la generalización de los resultados a otras instituciones o poblaciones.

    \item \textbf{Fatiga de alertas:} En el modelo final \texttt{random\_forest}, la especificidad de 0.3567 en el umbral operativo implica una tasa considerable de falsas alarmas. En entorno real, esto puede deteriorar la adherencia del personal a las alertas.

    \item \textbf{Balance entre desempeño e interpretabilidad:} El modelo de mejor AUC (\texttt{random\_forest}) presenta menor transparencia individual de predicciones respecto al árbol de decisión. Sin embargo, la brecha de discriminación es mínima (0.0050), permitiendo priorizar interpretabilidad sin sacrificar sustancialmente el desempeño.
\end{itemize}

\subsection{Trabajo futuro}

\begin{enumerate}
    \item \textbf{Incorporación de variables temporales:} Integrar señales vitales continuas (frecuencia cardíaca, presión arterial, saturación de oxígeno) para capturar tendencias y patrones de deterioro hemodinámico.

    \item \textbf{Modelos de series temporales:} Dado que el shock hemorrágico es un proceso evolutivo y no un evento estático, explorar arquitecturas como LSTM (Long Short-Term Memory) o Transformers permitiría modelar la evolución temporal del estado del paciente durante y después de la cirugía.

    \item \textbf{Optimización de umbral dinámico:} Investigar umbrales adaptativos que ajusten la sensibilidad según el contexto clínico (tipo de cirugía, perfil de riesgo del paciente) para reducir la fatiga de alertas sin comprometer la detección de casos críticos.

    \item \textbf{Validación prospectiva:} Implementar un estudio prospectivo para evaluar el rendimiento del modelo en condiciones clínicas reales y medir el impacto en los tiempos de respuesta del equipo médico.

    \item \textbf{Integración con sistemas hospitalarios:} Desarrollar la interfaz de integración con los sistemas de información quirúrgica de la FVL para alertas en tiempo real, incluyendo capacitación al personal sobre cómo interpretar y actuar ante las alertas.

    \item \textbf{Ampliación del dataset:} Colaborar con otras instituciones para aumentar el tamaño y diversidad del conjunto de datos, mejorando la generalización del modelo y permitiendo técnicas más avanzadas.
\end{enumerate}

\newpage

% -----------------------------------------------------------------------------
% ANEXOS
% -----------------------------------------------------------------------------
